% !TEX root = ../main.tex
\chapter{Symmetry}\label{ch:c11}
\begin{paracol}{2}

\begin{sources}
\begin{tabularx}{\linewidth}{@{}l l L@{}}
\toprule
\textbf{Source} & \textbf{Length} & \textbf{Content}\\
\midrule
\seg{11}{1}{Lecture 11: Symmetry} & 1:43:30 & Why the round sphere is symmetric and the potato is not; the push-forward of vectors and the pull-back of covectors (components, curves, level lines of a function); the induced metric on an embedded manifold (round sphere and ellipsoid); integral curves, complete vector fields, singularities, and the flow of a vector field; the Lie algebra of vector fields, Lie subalgebras and structure constants; the symmetry of a metric under a Lie subalgebra, isometries; the Lie derivative (axioms, components, use for symmetries). Student questions on adjoints, completeness, the Lie derivative and cosmological symmetries.\\
\seg{11}{2}{Tutorial 11: Symmetries} & 0:59:57 & Solutions of the sheet by a tutor (not Schuller): true/false on push-forward, pull-back, flows and the Lie derivative; the pull-back of a gradient; the components of the push-forward and of the pulled-back metric; the metric of an ellipsoid; the rotation algebra on $S^2$ and its structure constants; the integral curve and the flow of $X_3$ and when it is a symmetry.\\
\bottomrule
\end{tabularx}

\smallskip
\textbf{Transcripts} \file{materials/transcripts/C11-V01.txt} and \file{C11-V02.txt} (YouTube captions).
\textbf{Sheet 11:} see \url{https://tales.mbivert.com/on-heraeus-winter-school-tutorials/}.
\end{sources}

\begin{reading}[Where this chapter sits in the other notes]
\emph{Geometric Anatomy of Theoretical Physics} notes: Lie groups and Lie algebras are chapters 13 to 18 there, and the symmetries of a metric (flows, isometries, Lie derivative) are chapter 28 there [exact titles not checked]. Here the tangent bundle and the commutator of vector fields are chapter~\ref{ch:c06}, the metric and the round sphere are chapter~\ref{ch:c10}, and the covariant derivative that the Lie derivative is compared to is chapter~\ref{ch:c07}. Symmetric solutions of the Einstein equations (spherical symmetry, cosmology) are used from chapter~\ref{ch:c13} onwards.
\end{reading}

% ==================================================================
\section{Symmetry of a metric: the idea}\label{sec:c11-idea}

Good morning and welcome to the eleventh lecture \ts{11}{1}{0:00}. The topic is symmetry, but the lecture also collects a number of elementary techniques of differential geometry that will be needed urgently for Einstein's theory, and several loose ends that were left out so far. They are motivated by intuition. Recall the round sphere: first of all a smooth manifold, which at that stage has no shape; as soon as one provides the specific metric of chapter~\ref{ch:c10}, the \emph{round} metric $g^{\mathrm{round}}$, it has properties that follow the intuition of shape. One has the feeling that the round sphere has rotational symmetry: turn it around as one wishes and after any such transformation it looks the same. The \emph{potato}, which is also $S^2$ as a smooth manifold but with a different metric (``it depends on where you buy your potatoes''), does not feature that symmetry \ts{11}{1}{0:20}\ts{11}{1}{1:44}.

This is quite different from how rotational symmetry was seen before: as the group $SO(3)$, ``at whatever level of sophistication'', where one first has an inner product on a vector space and then classifies the linear maps $A$ that respect it, $\langle Ax,Ay\rangle=\langle x,y\rangle$. Here there is, up to now, no metric at all, and once there is one the thing that has symmetry is the \emph{metric itself}: the only thing that distinguishes sphere and potato is the metric, so symmetry is a property of the metric. A metric provides an inner product in each tangent space, but that is not the point: there are many inner products in many tangent spaces which could redefine $SO(3)$ tangent space by tangent space. What is at stake is the \emph{distribution} of these inner products over the sphere, and whether this distribution is in some sense rotationally invariant or not \ts{11}{1}{2:04}\ts{11}{1}{3:26}.

The question is how to describe the symmetries of a metric, and this is not academic: it is of the utmost importance when one tries to find solutions of Einstein's equations, because nobody can solve them without additional assumptions, for instance on the symmetry of the solution. In many physical situations one knows or postulates such symmetries, and then one makes progress \ts{11}{1}{4:07}\ts{11}{1}{4:48}.

\section{The push-forward}\label{sec:c11-push}

\begin{definition}[New concepts: push-forward]\label{def:c11-push}
Let $M,N$ be smooth manifolds (no metric yet) and $\phi\colon M\to N$ a smooth map. The \textbf{push-forward} of $\phi$ is the map
\[
  \phi_*\colon TM\to TN,
\]
taking a vector $X\in T_pM$ to the vector $\phi_*X\in T_{\phi(p)}N$ defined by its action on any smooth function $f\in C^\infty(N)$,
\[
  (\phi_*X)(f):=X(f\circ\phi).
\]
\end{definition}
\ts{11}{1}{5:09}\ts{11}{1}{6:54}
This is a very natural construction: a vector on $N$ can act on functions on $N$, and one declares it to act like the original vector would have acted on $f$ after $\phi$ (``nothing really happened''). If one draws the bundles over their base manifolds, with the projections down to $M$ and $N$ and the map $\phi$ below, then $\phi_*$ is the map induced on top: an entire fibre $T_pM$ is sent into the fibre over $\phi(p)$, because the base point of the vector is mapped to the base point $\phi(p)$ \ts{11}{1}{8:16}\ts{11}{1}{8:38}. The mnemonic: \emph{vectors are pushed forward} \ts{11}{1}{9:19}.

\paragraph{Components.} Choose a chart $(U,x)$ of $M$ and a chart $(V,y)$ of $N$ (the manifolds may be totally different, even of different dimension). Apply $\phi_*$ to the basis vector $\partial/\partial x^i|_p$ and extract the $a$-th component with the dual basis $\mathrm dy^a$:
\[
  (\phi_*)^a{}_i:=\mathrm dy^a\Bigl(\phi_*\frac{\partial}{\partial x^i}\Big|_p\Bigr)=\frac{\partial}{\partial x^i}\Big|_p\bigl(y^a\circ\phi\bigr)=\frac{\partial\hat\phi^a}{\partial x^i}\Big|_p,
\]
where in the first step the definition of $\mathrm dy^a$ was used (a covector acts on a vector like the vector acts on the function), in the second the definition of $\phi_*$, and $\hat\phi:=y\circ\phi\circ x^{-1}$ is the chart representation of $\phi$ (the map from $x(U)\subseteq\R^{\dim M}$ to $y(V)\subseteq\R^{\dim N}$). Here $i$ runs from $1$ to $\dim M$ and $a$ from $1$ to $\dim N$: these are \emph{not} tensor components, because the two indices belong to different spaces (even different manifolds), but they are the components of the induced linear map between the tangent spaces, and very easy to compute in practice \ts{11}{1}{10:01}\ts{11}{1}{11:04}\ts{11}{1}{12:07}\ts{11}{1}{13:49}.

\begin{theorem}[The push-forward of a tangent vector to a curve]\label{thm:c11-curve}
Let $\gamma$ be a smooth curve in $M$ through $p=\gamma(\lambda_0)$ and $v_\gamma$ its tangent vector. Then
\[
  \phi_*\bigl(v_{\gamma,p}\bigr)=v_{\phi\circ\gamma,\phi(p)},
\]
the tangent vector of the pushed-forward curve $\phi\circ\gamma$ at $\phi(p)$: the push-forward sends tangent vectors of curves to tangent vectors of the pushed-forward curves.
\end{theorem}
\ts{11}{1}{14:31}\ts{11}{1}{15:33}\ts{11}{1}{16:56}
\begin{proof}
Let $f\in C^\infty(N)$. By the definition of the push-forward, $(\phi_*v_{\gamma,p})(f)=v_{\gamma,p}(f\circ\phi)$, and by the definition of the tangent vector of a curve this is the derivative $(f\circ\phi\circ\gamma)'(\lambda_0)$. By associativity of the composition this is $\bigl(f\circ(\phi\circ\gamma)\bigr)'(\lambda_0)=v_{\phi\circ\gamma,\phi(p)}(f)$. Since this holds for every $f$, the two vectors coincide.
\end{proof}
\ts{11}{1}{20:50}\ts{11}{1}{21:31}\ts{11}{1}{22:53}
(Very simple and, as Schuller says, very powerful, because the geometric picture is clear.)

\begin{coursenote}[Student question: different dimensions, embeddings]
A student asks about the case where $M$ and $N$ have different dimensions, where the curve can end up with a ``different dimension'' too. The interesting case is $\dim M<\dim N$ with $\phi$ injective (for instance $\mathbb R^2$ put into a bigger space like a hemisphere): a tangent vector of $M$ can be considered an \emph{intrinsic} tangent vector (``the velocity the ant would talk about, who knows nothing about the third dimension'') or an \emph{extrinsic} one, a vector of the bigger space (``the velocity you would talk about seeing the ant running along in three space''): it is twice the same object from two perspectives, and the push-forward converts the first into the second. This holds for any smooth injective $\phi$, which is what an embedding is \ts{11}{1}{17:37}\ts{11}{1}{18:02}\ts{11}{1}{19:48}\ts{11}{1}{20:08}.
\end{coursenote}
\ts{11}{1}{20:30}
Another application will be the case $N=M$, taken up with the flow below.

\section{The pull-back}\label{sec:c11-pull}

\begin{definition}[New concepts: pull-back]\label{def:c11-pull}
For a smooth map $\phi\colon M\to N$ the \textbf{pull-back} is the map
\[
  \phi^*\colon T^*N\to T^*M,
\]
taking a covector $\omega\in T^*_{\phi(p)}N$ to the covector $\phi^*\omega\in T^*_pM$ defined on a vector $X\in T_pM$ by
\[
  (\phi^*\omega)(X):=\omega\bigl(\phi_*X\bigr).
\]
\end{definition}
\ts{11}{1}{23:14}\ts{11}{1}{24:17}\ts{11}{1}{25:20}
It goes the opposite way: its domain is the cotangent bundle of the \emph{target}, hence the name pull-back; and it is defined \emph{through the push-forward} (careful: ``index up is the pull-back, index down is the push-forward''). Mnemonic: \emph{covectors are pulled back}. (Schuller first mixed up on the board which one takes what, and corrected himself.)

\paragraph{Components.} With the same charts, pull back the chart-induced covector $\mathrm dy^a$ at $\phi(p)$ and evaluate it on $\partial/\partial x^i$:
\[
  (\phi^*)_i{}^a=\bigl(\phi^*\mathrm dy^a\bigr)\Bigl(\frac{\partial}{\partial x^i}\Big|_p\Bigr)=\mathrm dy^a\Bigl(\phi_*\frac{\partial}{\partial x^i}\Big|_p\Bigr)=(\phi_*)^a{}_i=\frac{\partial\hat\phi^a}{\partial x^i}\Big|_p.
\]
The components of the pull-back are numerically \emph{exactly the same} as those of the push-forward: one computes them once, and whether one contracts with one index or with the other decides whether one performs a push-forward of a vector (component $a$, landing in $N$) or a pull-back of a covector (component $i$, landing in $M$) \ts{11}{1}{26:03}\ts{11}{1}{27:25}\ts{11}{1}{28:07}\ts{11}{1}{29:09}.

\paragraph{The picture.} Take a function $f$ on $N$ and its level lines (``like the isotherms of a temperature''), and a map $\phi$ that spreads them apart. The covector field $\mathrm df$ (the gradient) on $N$ can be pulled back, and the result is the gradient of the function $f\circ\phi$, whose level lines are the pulled-back ones; this is shown in the tutorial (\cref{sec:t11-1}, exercise 2) \ts{11}{1}{29:52}\ts{11}{1}{30:33}.

\begin{coursenote}[Student questions: are push-forward and pull-back inverse? adjoint?]
Does it make sense to ask whether push-forward and pull-back are somehow inverse to each other? Not really: $\phi_*$ goes from $TM$ to $TN$ and $\phi^*$ from $T^*N$ to $T^*M$, and none of these four spaces is the same, so there is nothing to invert. A student suggests that the pull-back is the \emph{adjoint}, or dual, map of the push-forward (``if you have $V\to W$ you can do the dual map''): yes, that is exactly the linear algebra analogue, but dual is not inverse \ts{11}{1}{31:01}\ts{11}{1}{32:11}.
Generically one cannot pull a vector back nor push a covector forward with these data alone. If $\phi$ is invertible one can: a covector on $M$ is brought to $N$ by pulling it back along $\phi^{-1}$. This was the point of the question what happens if $M$ is smaller than $N$: then $\phi$ cannot be invertible \ts{11}{1}{32:52}\ts{11}{1}{33:32}.
\end{coursenote}

\section{The induced metric}\label{sec:c11-induced}

An important application of push-forward and pull-back: let $\phi\colon M\to N$ be an \emph{embedding}, an injective smooth map with $\dim M<\dim N$ (so that $M$ sits inside $N$), and let a metric $g$ live on the bigger manifold $N$. For instance $N=\R^3$ with its standard topology and smooth atlas, $M=S^2$: there are many maps $\phi\colon S^2\to\R^3$, and it depends on $\phi$ how the sphere lies in $\R^3$ (like a round sphere, or like a potato; still, the smooth manifold $S^2$ has no shape of its own, only the way it lies in there). Can one take the Euclidean metric of $\R^3$ (components $1,1,1$ in the identity chart) and inherit it to the surface? A different embedding will probably give a different metric. Which one does one want? The picture: take a string and stretch it, so that it has a length in $\R^3$, and put it on the sphere; one wants a metric such that the length computed on the sphere is the length of the string \ts{11}{1}{34:42}\ts{11}{1}{35:23}\ts{11}{1}{36:27}\ts{11}{1}{37:07}.

\begin{definition}[New concepts: induced metric]\label{def:c11-induced}
Let $\phi\colon M\to N$ be an embedding and $g_N$ a metric on $N$. The \textbf{induced metric} $g_M$ on $M$ is the pull-back $g_M:=\phi^*g_N$ of the metric, a $(0,2)$ tensor field on $M$ acting on vectors $X,Y$ in the same tangent space $T_pM$ as
\[
  g_M(X,Y):=g_N\bigl(\phi_*X,\phi_*Y\bigr)\big|_{\phi(p)},
\]
i.e.\ the vectors are translated to $N$ by the push-forward and fed into the metric there.
\end{definition}
\ts{11}{1}{37:28}\ts{11}{1}{38:10}
So the embedding matters for the result. In components, with charts as above (indices $i,j$ on $M$, $a,b$ on $N$),
\[
  (g_M)_{ij}(p)=\frac{\partial\hat\phi^a}{\partial x^i}\Big|_p\,\frac{\partial\hat\phi^b}{\partial x^j}\Big|_p\,(g_N)_{ab}\big|_{\phi(p)},
\]
where the metric components are evaluated at the point $\phi(p)$, not at $p$ (Schuller wrote the point wrong on the board and corrected it): twice the components of the push-forward \ts{11}{1}{38:31}\ts{11}{1}{38:51}\ts{11}{1}{39:13}.
In the tutorials one shows that the standard embedding of the sphere into $\R^3$ gives the round metric of chapter~\ref{ch:c10}, which was written down without saying where it comes from; and that an embedding with different axes gives an ellipsoid metric, also a metric on the same smooth $S^2$ (``the ellipsoid we could not call an ellipsoid before we induced a metric''). Hence \emph{the shape lies in the metric}, and this is how one gets new metrics on embedded submanifolds from a metric outside \ts{11}{1}{39:33}\ts{11}{1}{39:54}\ts{11}{1}{40:14}\ts{11}{1}{40:35}.

\begin{coursenote}[Student questions: the inverse metric, and the string]
Can one do the same with the inverse metric? One could try to define it through covectors pulled back from the covectors over there, but point by point: the push-forward and pull-back act on the whole tangent bundle, at any point one can send any vector back and forth. Mapping \emph{fields} to fields, so that a section remains a section, is something different \ts{11}{1}{41:01}\ts{11}{1}{41:24}\ts{11}{1}{41:45}.
How does one know along which path the string has maximal elongation? The ``quick'' picture is: curves in the ambient space are measured by their tangent vectors; a curve on the sphere is pushed forward by the embedding to a curve in $\R^3$, its tangent vectors are the push-forwards of the original ones (\cref{thm:c11-curve}), and they are measured with the metric of $\R^3$. By this way of thinking about the length of curves, this is exactly the induced metric \ts{11}{1}{42:05}\ts{11}{1}{42:26}\ts{11}{1}{42:46}\ts{11}{1}{43:07}.
\end{coursenote}

\section{Integral curves and the flow of a vector field}\label{sec:c11-flow}

Now closer to symmetry. Let $M$ be a smooth manifold and $X$ a vector field on it \ts{11}{1}{43:50}.

\begin{definition}[New concepts: integral curve]\label{def:c11-integral}
A curve $\gamma\colon I\to M$, with $I\subseteq\R$ an interval, is an \textbf{integral curve} of the vector field $X$ if at each point of the curve its tangent vector reproduces the given vector of the field,
\[
  v_{\gamma,\gamma(\lambda)}=X_{\gamma(\lambda)}\qquad\text{for all }\lambda\in I.
\]
\end{definition}
\ts{11}{1}{44:10}\ts{11}{1}{45:17}
Think of the vector field as the velocity of the water molecules in a river: an integral curve is the trajectory of a (paper) ship thrown in that has no propulsion and just flows with it \ts{11}{1}{45:39}\ts{11}{1}{46:00}.

\begin{definition}[New concepts: complete vector field]\label{def:c11-complete}
A vector field $X$ is \textbf{complete} if all its integral curves have as their domain $I$ the entire real line $\R$.
\end{definition}
\ts{11}{1}{46:44}
\begin{example}[A field made incomplete]\label{ex:c11-incomplete}
Take a smooth vector field on a plane, and an integral curve through some point: it can go around and around and come from minus infinity, its parameter can be extended arbitrarily. A point where the field vanishes is a constant integral curve, again with domain $\R$. So this field is complete, \emph{up to the moment one viciously removes a point from the space}: then an integral curve running towards it comes arbitrarily close to the missing point but never reaches it, so it stops at a finite parameter value, and the same vector field on the space with the point removed is incomplete. It is not the field that misses a vector (then it would not be a vector field): the point has to be removed from the space.
\end{example}
\ts{11}{1}{47:06}\ts{11}{1}{48:15}\ts{11}{1}{48:56}
\begin{aside}[Singularities]
In spacetime physics there are many such points which do not belong to the spacetime manifold, the so-called \emph{singularities}, for instance where the curvature blows up. The Hawking--Penrose singularity theorems were proven for singularities defined in this or a similar manner (depending on what one wants to do with geodesics): singularities are points where curves stop at a finite parameter value, and completeness is the technical door to the singularity theorems. ``Very roughly, I am talking of the zeroth order contribution to the definition''; the completeness of vector fields is a very important point \ts{11}{1}{48:56}\ts{11}{1}{49:17}\ts{11}{1}{49:38}.
\end{aside}
\begin{coursenote}[Student questions: restricting or reparametrising the domain; compactness]
Cannot one restrict the domain of an integral curve, or reparametrise to get all of $\R$? Restricting leaves it without the domain $\R$; and one cannot reparametrise, because the given field forces the parametrisation that reproduces \emph{exactly} that velocity, not double or half as long (one cannot ``go there very slowly'')
\ts{11}{1}{49:58}\ts{11}{1}{50:18}.
A theorem: a smooth vector field with compact support (in particular, any smooth vector field on a compact manifold) is complete. The manifold above with a point taken out is not compact; compactness is a topological notion (every open cover has a finite subcover, ``whatever that means'') \ts{11}{1}{50:40}\ts{11}{1}{51:04}.
\end{coursenote}

\begin{definition}[New concepts: flow of a complete vector field]\label{def:c11-flow}
The \textbf{flow} of a complete vector field $X$ is the one-parameter family of maps (``$h$ like flow, I do not know; it is the traditional letter'')
\[
  h^X\colon\R\times M\to M,\qquad(\lambda,p)\mapsto h^X(\lambda,p)=:h^X_\lambda(p):=\gamma_p(\lambda),
\]
where $\gamma_p\colon\R\to M$ is the integral curve of $X$ with $\gamma_p(0)=p$ (one can always choose that).
\end{definition}
\ts{11}{1}{51:24}\ts{11}{1}{52:30}\ts{11}{1}{53:15}\ts{11}{1}{54:18}
The map takes an arbitrary point, constructs the integral curve starting at it for $\lambda=0$, and after the parameter has passed an amount $\lambda$ one has gone from $p$ to $\gamma_p(\lambda)$. For fixed $\lambda$, $h^X_\lambda\colon M\to M$ is a smooth map that takes every point of the manifold and pushes it further by a parameter distance $\lambda$ along the integral curves of the field \ts{11}{1}{54:38}\ts{11}{1}{54:58}.
One does not push forward vectors in this case but \emph{points}: this is a special instance of the map $\phi$ of the previous sections, a map from a manifold into itself. In general the image of a set is not the same set: ``you take the river, you take a set, you wait for a parameter time $\lambda$, and the set has evolved'', for instance if the river flows apart. ``This is Heraclitus' theorem, if you wish, if the river is flowing, if the vector field is not zero'' (as said in class: ``you never jump in the same river twice''; the attribution to Heraclitus is Schuller's, not checked) \ts{11}{1}{55:20}\ts{11}{1}{55:42}\ts{11}{1}{56:03}\ts{11}{1}{56:24}.

This takes us closer to what happens on the sphere. So far no metric was used in this flow business: only a smooth manifold and a given vector field. Think of the sphere as the potato, with no metric, and give vector fields on it, ``around all possible axes (but there are no axes)''. Let the points flow; one can look at the induced push-forward of vectors, the pull-back of covectors, the induced metric: if the metric has not changed, the metric was symmetric under it. These flows are the key to understand what happens if one moves things; but how does one distinguish spherical from cylindrical from translational symmetry? \ts{11}{1}{56:44}\ts{11}{1}{57:06}\ts{11}{1}{57:46}\ts{11}{1}{58:07}\ts{11}{1}{58:28}

\section{The Lie algebra of vector fields}\label{sec:c11-lie}

The set $\Gamma(TM)$ of all vector fields is a $C^\infty(M)$-module (chapter~\ref{ch:c06}). Restrict to considering it as an $\R$-vector space: then in general there is no finite basis, one needs many vector fields to produce all the others \ts{11}{1}{59:09}\ts{11}{1}{59:30}\ts{11}{1}{59:52}. The point is that for two vector fields $X,Y$ one can form the commutator, again in $\Gamma(TM)$,
\[
  [X,Y](f):=X\bigl(Y(f)\bigr)-Y\bigl(X(f)\bigr),
\]
which is a vector field because it acts on functions. It is
\begin{enumerate}
  \item antisymmetric: $[X,Y]=-[Y,X]$;
  \item $\R$-linear in each slot: $[\lambda X+Z,Y]=\lambda[X,Y]+[Z,Y]$ for $\lambda\in\R$ (only for real numbers);
  \item it satisfies the \textbf{Jacobi identity}: $[[X,Y],Z]+[[Y,Z],X]+[[Z,X],Y]=0$ (the cyclic rotation of what is inside the brackets).
\end{enumerate}
\ts{11}{1}{1:00:13}\ts{11}{1}{1:00:34}\ts{11}{1}{1:00:56}

\begin{definition}[New concepts: Lie algebra]\label{def:c11-liealg}
A \textbf{Lie algebra} is a vector space (also any module over a commutative ring) with a map, the \emph{bracket}, that takes two elements to another one of the same space, which is antisymmetric, linear in each slot and satisfies the Jacobi identity. (The name is after Sophus Lie, ``Lie'' pronounced ``lee''.)
\end{definition}
\ts{11}{1}{1:01:18}\ts{11}{1}{1:01:39}\ts{11}{1}{1:02:00}
Obviously all the vector fields on a manifold, with the commutator, constitute a ``massively huge'' Lie algebra, infinite-dimensional if one regards it as an $\R$-vector space.

\begin{definition}[New concepts: Lie subalgebra, structure constants]\label{def:c11-subalg}
Let $X_1,\dots,X_s$ be $s$ vector fields on $M$ such that the commutator of any two is again a \emph{real} linear combination of them,
\[
  [X_i,X_j]=c^k{}_{ij}\,X_k,\qquad i,j,k=1,\dots,s,\qquad c^k{}_{ij}\in\R,
\]
numbers, not functions. The $c^k{}_{ij}$ are the \textbf{structure constants}. The \emph{real} span $L=\mathrm{span}_\R\{X_1,\dots,X_s\}$ is then a sub vector space of $\Gamma(TM)$ on which the inherited bracket does not lead outside $L$: an $s$-dimensional \textbf{Lie subalgebra} (``it would be more logical to call it a subLie algebra, but language demands Lie subalgebra'').
\end{definition}
\ts{11}{1}{1:02:20}\ts{11}{1}{1:03:02}\ts{11}{1}{1:03:43}\ts{11}{1}{1:04:03}
There are many such choices of subalgebras of the huge algebra.

\begin{example}[The rotation algebra on $S^2$]\label{ex:c11-so3}
On $S^2$ take three vector fields $X_1,X_2,X_3$ with commutation relations of the angular momentum algebra of quantum mechanics (``the little $i$ is missing because we are over the real numbers; insert it and you get the minuses''): this algebra, the span of the three, is usually called $\mathfrak{so}(3)$. There is no metric anywhere. In the chart $(\theta,\varphi)$ of the sphere, evaluated at a point $p$,
\begin{align*}
  X_1&=-\sin\varphi\,\frac{\partial}{\partial\theta}-\cot\theta\cos\varphi\,\frac{\partial}{\partial\varphi},\\
  X_2&=\cos\varphi\,\frac{\partial}{\partial\theta}-\cot\theta\sin\varphi\,\frac{\partial}{\partial\varphi},\\
  X_3&=\frac{\partial}{\partial\varphi},
\end{align*}
where $\theta,\varphi$ stand for the chart coordinate functions evaluated at $p$ (``instead of writing a chart I write them more intuitively as maps $\theta$ and $\varphi$''). One can verify by plugging in that they satisfy the commutation relations (done partly in the tutorial, \cref{sec:t11-1}). On the merely smooth manifold $S^2$, with no additional structure, one thus has something one would immediately call an $SO(3)$; the flows of these fields take points like rotations, ``like this or like that''. Which algebra one has is recognised from the structure constants in certain bases.
\end{example}
\ts{11}{1}{1:04:47}\ts{11}{1}{1:05:15}\ts{11}{1}{1:06:17}\ts{11}{1}{1:07:01}\ts{11}{1}{1:07:42}\ts{11}{1}{1:08:02}\ts{11}{1}{1:08:22}
\begin{coursenote}[Slips at the board and the sign of the structure constants]
Schuller starts with ``minus cosine theta [?]'' and corrects on the spot to the $X_1$ above (``I have to quickly look; I think that is right''). The commutation relations themselves were not written out in the captions (``these relations'' on the board); computing from the three fields as given one finds (not said in class) $[X_1,X_2]=-X_3$, $[X_2,X_3]=-X_1$, $[X_3,X_1]=-X_2$, i.e.\ $c^k{}_{ij}=-\epsilon_{ijk}$, which is $\mathfrak{so}(3)$ up to the sign convention of the generators. The tutorial gets $[X_1,X_3]=X_2$, consistent with this.
\end{coursenote}

\section{The symmetry of a metric}\label{sec:c11-symmetry}

Now the definition of the symmetry of a metric.

\begin{definition}[New concepts: symmetry of a metric, isometry]\label{def:c11-symmetry}
An $s$-dimensional Lie subalgebra $L\subseteq\Gamma(TM)$ is said to be a \textbf{symmetry} of a metric tensor field $g$ if for every vector field $X\in L$ (not only $X_1,X_2,X_3$, but anything in their \emph{real} span) and for all $\lambda\in\R$,
\[
  \bigl(h^X_\lambda\bigr)^*g=g,\qquad\text{i.e.}\qquad g\bigl(h^X_{\lambda*}A,\;h^X_{\lambda*}B\bigr)=g(A,B)\quad\text{for all vectors }A,B\text{ at the same point}.
\]
Each map $h^X_\lambda$ is then called an \textbf{isometry}: one has a one-parameter family of isometries generated from each vector field in $L$.
\end{definition}
\ts{11}{1}{1:08:42}\ts{11}{1}{1:09:28}\ts{11}{1}{1:09:49}\ts{11}{1}{1:10:10}\ts{11}{1}{1:12:39}\ts{11}{1}{1:18:31}
The pull-back of the metric by the flow is defined as in \cref{def:c11-induced} with $M=N$, by acting on vectors: $(\phi^*g)(A,B)=g(\phi_*A,\phi_*B)$. So one flows along the field, looks at the metric elsewhere, and if the result does not differ from how it looked before, then apparently the metric has not changed, and the algebra describes a symmetry of the metric \ts{11}{1}{1:11:37}\ts{11}{1}{1:12:18}\ts{11}{1}{1:13:00}.

\begin{coursenote}[Why the vector fields are complete; the dimension $s$; the board slip]
For the flow one needs the fields in $L$ to be complete. ``The dirty truth, not the dirty truth, the beautiful truth'': the Lie algebra belongs to a Lie group, itself a manifold with a group structure; the Lie algebra is the set of left-invariant vector fields, which can be identified with the tangent space at the identity; every left-invariant vector field on a Lie group is complete; acting with the group on a manifold gives vector fields on the manifold in which the algebra is mirrored. Schuller does not show this (it cannot be understood with the material so far), but the guarantee is that these fields are complete \ts{11}{1}{1:10:36}\ts{11}{1}{1:10:56}\ts{11}{1}{1:11:16}.
A student notes that the dimension $s$ does not enter the definition: right, $s$ is just that of the finite-dimensional subalgebra chosen \ts{11}{1}{1:13:20}\ts{11}{1}{1:13:41}.
At first Schuller wrote $X$ as an argument of the pulled-back metric, then realised ``the $X$ is the one from the algebra and here you act on vector fields $A$ and $B$'': the equation above is the right one, with the push-forward of the arguments, evaluated on two vectors at the same point \ts{11}{1}{1:16:06}\ts{11}{1}{1:17:08}\ts{11}{1}{1:17:50}.
\end{coursenote}
Since the chart $(\theta,\varphi)$ covers almost the entire sphere (a set is taken out), a proper treatment provides the same vector fields in a second chart, compatible on the overlap. With this set of vector fields one can check rotational symmetry or not of any given metric \ts{11}{1}{1:14:01}.

\begin{example}[Round sphere, ellipsoid, potato]\label{ex:c11-ellipsoid}
A metric that is $SO(3)$ symmetric in the sense of \cref{def:c11-symmetry} is the round one. What about an ellipsoid whose axes are $a=b$ and $c$ different (``elongated like this, but seen from above a circle'')? Intuitively it has symmetry around one axis only (``cylindrical''), which reduces the algebra to the one field $X_3$ (if chosen cleverly): a one-dimensional Lie subalgebra, trivial since all structure constants vanish. So the ellipsoid with $a=b\neq c$ has only that symmetry, and is not $SO(3)$ symmetric. The tutorial checks this for $X_3$ (\cref{sec:t11-1}).
\end{example}
\ts{11}{1}{1:14:22}\ts{11}{1}{1:15:03}\ts{11}{1}{1:15:24}\ts{11}{1}{1:19:12}
\begin{aside}[Why this matters: symmetric ans\"atze]
It looks like a lot of work for small statements, but it is of direct application to spacetime physics: later the metric of a Lorentzian spacetime will be given, and the question what a \emph{spherically symmetric} situation is will be answered using exactly these fields applied to the four-dimensional metric. ``Textbooks often in general relativity say: let us make the following ansatz for a spherically symmetric metric; then follows a lot of handwaving. No handwaving on earth can replace this'': one has to say which symmetry algebra one wants implemented \ts{11}{1}{1:20:35}\ts{11}{1}{1:20:56}\ts{11}{1}{1:37:50}.
\end{aside}

\section{The Lie derivative}\label{sec:c11-lieder}

The condition of \cref{def:c11-symmetry} is geometrically intuitive but a mess to check: one has to check it for all $\lambda$ and all vector fields. The Lie derivative makes checking for symmetries child's play \ts{11}{1}{1:21:19}\ts{11}{1}{1:21:39}.

If $X\in L$ generates a symmetry, then $(h^X_\lambda)^*g-g=0$ for all $\lambda$; dividing by $\lambda$ and taking the limit $\lambda\to0$ it stays zero. Write
\[
  \mathcal L_Xg:=\lim_{\lambda\to0}\frac{(h^X_\lambda)^*g-g}{\lambda}.
\]
So $\mathcal L_Xg=0$ for all $X\in L$ if $L$ is a symmetry (and conversely: if the derivative vanishes, $g$ does not change along the flow, ``you just integrate''). This does not help much by itself, one has added insult to injury, but it can be shown that this is precisely the so-called \emph{Lie derivative} of the metric with respect to the vector field $X$ \ts{11}{1}{1:22:02}\ts{11}{1}{1:22:45}\ts{11}{1}{1:23:30}\ts{11}{1}{1:24:11}.

\begin{definition}[New concepts: Lie derivative]\label{def:c11-liederivative}
On a smooth manifold $(M,\mathcal O,\mathcal A)$ (no extra structure is needed, unlike for the covariant derivative) \emph{the} Lie derivative $\mathcal L$ takes a vector field $X$ and a $(p,q)$ tensor field $T$ to a $(p,q)$ tensor field $\mathcal L_XT$, such that
\begin{enumerate}
  \item on functions, $\mathcal L_Xf=X(f)$;
  \item on vector fields, $\mathcal L_XY=[X,Y]$;
  \item it is additive: $\mathcal L_X(T+S)=\mathcal L_XT+\mathcal L_XS$ for tensors of the same valence;
  \item it satisfies the Leibniz rule, written for a $(1,1)$ tensor $T$ and extended to any valence: $\mathcal L_X\bigl(T(\omega,Y)\bigr)=(\mathcal L_XT)(\omega,Y)+T(\mathcal L_X\omega,Y)+T(\omega,\mathcal L_XY)$;
  \item it is additive in the lower slot: $\mathcal L_{X+Z}T=\mathcal L_XT+\mathcal L_ZT$.
\end{enumerate}
\end{definition}
\ts{11}{1}{1:24:52}\ts{11}{1}{1:26:15}\ts{11}{1}{1:26:37}\ts{11}{1}{1:26:58}\ts{11}{1}{1:27:42}\ts{11}{1}{1:28:03}
(Schuller says ``a'' Lie derivative and corrects: it is \emph{the} Lie derivative, there is only one such.) This looks like the covariant derivative axioms, and one worries about the remaining freedom: the requirement in the lower slot is even \emph{weaker} (only additivity: it is not required that a smooth function jumps out), so potentially even more freedom than for the covariant derivative, which needed extra structure to fix it. It is axiom 2 that fixes everything, without additional structure \ts{11}{1}{1:28:27}\ts{11}{1}{1:28:49}\ts{11}{1}{1:29:10}.

\begin{coursenote}[Not $C^\infty$-linear in $X$, and why this is the price of no structure]
For the torsion one needed $C^\infty$-linearity and could show it, pulling the $f$ out only at the cost of an extra term; here one cannot require it, hence the weaker system. Why was this not defined before the covariant derivative? It is a derivative along $X$ that needs no additional structure, but it ``cheats'': it sucks in the information of how the vector field behaves \emph{in a neighbourhood} of the point (the flow pushes one forward), not only its value at the point. This is reflected in the missing $C^\infty$-linearity: changing $X$ by a function that is a number at the point but varies around it changes the result, through the extra term. This is not ``a worse derivative'' (``a nonsense statement''): the covariant derivative is $C^\infty(M)$-linear in the lower slot and the Lie derivative is not, and that is the difference. The covariant derivative was needed because one wanted a $C^\infty$-linear object; the Lie derivative came up from the considerations of symmetry and of things not changing along flows. The field $X$ need not be complete (only a vector field in a neighbourhood of the point of evaluation) \ts{11}{1}{1:29:30}\ts{11}{1}{1:30:10}\ts{11}{1}{1:30:31}\ts{11}{1}{1:31:12}\ts{11}{1}{1:31:33}\ts{11}{1}{1:31:54}.
\end{coursenote}

\paragraph{Components.} In a chart $(U,x)$ the result of acting on a vector field (a five-minute calculation from the axioms, which Schuller recommends as an exercise: ``some examiner might get the idea to set this as an examination task'') is
\[
  (\mathcal L_XY)^i=X^m\partial_mY^i-Y^m\partial_mX^i,
\]
to be compared with $(\nabla_XY)^i=X^m\partial_mY^i+\Gamma^i{}_{sm}X^mY^s$: there is a \emph{minus} instead of a plus, and no $\Gamma$ (``the $i$ came from there, the $m$ here and the $s$ there''): it needs the derivative of $X$, hence not the value at a point only. In general,
\[
  (\mathcal L_XT)^i{}_j=X^m\partial_mT^i{}_j-T^s{}_j\,\partial_sX^i+T^i{}_s\,\partial_jX^s,
\]
the opposite of the covariant derivative: a minus for every index up and a plus for every index down. The application is that for the metric
\[
  (\mathcal L_Xg)_{ij}=X^m\partial_mg_{ij}+g_{sj}\,\partial_iX^s+g_{is}\,\partial_jX^s,
\]
with two pluses, since both indices of $g$ are down. (Schuller does not prove that the limit definition reproduces this: ``you need to go to specially adapted charts and argue, I spare you that, you need to trust me anyway''.) It is thus very easy to check whether a metric features a symmetry, and conversely: take the most general metric, impose the symmetry, and learn how it restricts the components of the metric in that coordinate system, to be used as an ansatz in the Einstein equations with many components already zero \ts{11}{1}{1:32:15}\ts{11}{1}{1:32:36}\ts{11}{1}{1:32:57}\ts{11}{1}{1:34:40}\ts{11}{1}{1:35:01}\ts{11}{1}{1:36:02}\ts{11}{1}{1:36:44}\ts{11}{1}{1:37:25}.

\begin{coursenote}[Student questions at the end]
\begin{itemize}
  \item \emph{Is the equivalence obvious in one direction?} If one defines the symmetry as $\phi^*g=g$ for all flows, then the difference quotient is zero, so the Lie derivative is zero, and by integrating the other way round the converse holds; the minus is the subtraction in the vector space of $(0,2)$ tensors and the division the scalar multiplication there (not in the real numbers) \ts{11}{1}{1:38:51}\ts{11}{1}{1:39:11}\ts{11}{1}{1:39:32}\ts{11}{1}{1:39:53}.
  \item \emph{Derivation of the component formula?} Not derived: ``I wrote this from memory, so this should give you an indication that you need to know it by memory''. For vector fields it is the commutator; for covector fields and tensors one applies the Leibniz rule (once, twice) starting from the axioms \ts{11}{1}{1:40:14}\ts{11}{1}{1:40:34}\ts{11}{1}{1:40:54}.
  \item \emph{Does a given set of symmetries identify the algebra?} It is not always easy to read the symmetry off the structure constants (the translation algebra has all structure constants vanishing, and is no longer semisimple). The basic idea is: if physically one has an idea what the vector fields are, whether they commute or not, one implements them and requires invariance. In cosmology, in a four-dimensional spacetime, one will require vector fields in three directions plus further ones, giving a spacetime metric that is intuitively (there is no metric yet, so no ``spatial'') homogeneous and isotropic in three dimensions: a six-dimensional symmetry algebra, requiring that at large scales all directions look the same and no point of the universe is preferred. One can make big symmetry algebras smaller and use intuition for what one means; but the algebra is more fundamental than the group it came from \ts{11}{1}{1:41:26}\ts{11}{1}{1:41:47}\ts{11}{1}{1:42:07}\ts{11}{1}{1:42:28}\ts{11}{1}{1:42:48}\ts{11}{1}{1:43:08}.
\end{itemize}
\end{coursenote}
The rest of the lecture was announced as the integration on manifolds, ``the way Hilbert did it'' (chapter~\ref{ch:c12}), after a break \ts{11}{1}{1:38:10}\ts{11}{1}{1:38:31}.

% ==================================================================
\section{Tutorial 11}\label{sec:t11-1}

Tutorial 11, 0:59:57, \seg{11}{2}{video}, by a tutor, not Schuller. The sheet (Sheet 11, ``Symmetries'') is paraphrased, not copied; see \url{https://tales.mbivert.com/on-heraeus-winter-school-tutorials/} and \url{https://tales.mbivert.com/heraeus/Sheet_11.pdf}. (Mis-hearings in the captions: $\phi$ is transcribed as ``five'', ``phi'', ``fi'' or ``Chi''; $g$ as ``G'', ``T'' or ``Chi''; ``Lie'' as ``Lisa'' or ``lead''.)

\subsection*{Exercise 1: true or false}

\paragraph{(A) Push-forward and pull-back.}
\begin{itemize}
  \item The push-forward maps a vector field to a vector field: \textbf{false}. $\phi$ is only smooth, not injective: two points $p,q$ of $M$ may be mapped to the same point of $N$, giving two vectors in one tangent space, whereas a vector field has exactly one vector at each point \ts{11}{2}{0:20}\ts{11}{2}{1:02}\ts{11}{2}{1:45}.
  \item The pull-back of a covector yields a covector: \textbf{true}, $\phi^*$ maps $T^*N$ to $T^*M$, every element of $T^*_{\phi(p)}N$ to one of $T^*_pM$ \ts{11}{2}{2:28}\ts{11}{2}{2:48}.
  \item One can use a metric on a smoothly embedded manifold by pulling it back from the ambient manifold: \textbf{true}, $(\phi^*g)(X,Y)=g(\phi_*X,\phi_*Y)$; an embedding is an injective smooth map, and one always can find one \ts{11}{2}{3:09}\ts{11}{2}{3:30}\ts{11}{2}{4:11}.
  \item Push-forward and pull-back are linear maps: \textbf{true}. For $X,Y\in T_pM$ and $f\in C^\infty(N)$: $\phi_*(X+Y)(f)=(X+Y)(f\circ\phi)=X(f\circ\phi)+Y(f\circ\phi)=(\phi_*X)(f)+(\phi_*Y)(f)$ (the addition of vectors is that of the vector space $T_pM$, which is why they must lie in the same space), and likewise for scalars and for the pull-back \ts{11}{2}{4:32}\ts{11}{2}{4:54}\ts{11}{2}{6:40}.
  \item The push-forward is a $(1,1)$ tensor: \textbf{false}, a $(1,1)$ tensor is a multilinear map $T^*M\times TM\to\R$, which is not the definition of the push-forward \ts{11}{2}{7:01}.
  \item One can pull back a $(1,1)$ tensor only along a bijective map: \textbf{incorrect}, says the tutor, since one can actually pull back tensors along a bijective map, so the statement is wrong in the form on the sheet [the exact wording was not recovered from the captions, ?]. (Added, not said in class: any covariant tensor, e.g.\ a $(0,2)$ one, can be pulled back along any smooth map; pulling back a mixed tensor needs a diffeomorphism.) \ts{11}{2}{7:22}\ts{11}{2}{7:43}.
\end{itemize}

\paragraph{(B) Symmetries.}
\begin{itemize}
  \item A Lie algebra $L$ is a symmetry of a metric $g$ if the pull-back of $g$ along the flow of any element of $L$ reproduces $g$: \textbf{true}. The tutor recalls integral curves, the flow map $h^X_\lambda\colon M\to M$ (smooth, so one can push forward or pull back along it): if after going back some steps the metric looks the same, $X$ is a symmetry, and if this holds for all of $L$, $L$ is \ts{11}{2}{7:43}\ts{11}{2}{8:03}\ts{11}{2}{9:25}\ts{11}{2}{9:46}.
  \item The Lie derivative of $Y$ with respect to $X$ is the commutator $[X,Y]$: \textbf{true}, axiom 2 \ts{11}{2}{10:27}\ts{11}{2}{10:48}.
  \item $L$ is a symmetry of $g$ if $\mathcal L_Xg=0$ for every $X\in L$: \textbf{true}, since $\mathcal L_Xg=\lim_{\lambda\to0}\frac1\lambda\bigl((h^X_\lambda)^*g-g\bigr)$ and the numerator vanishes by the previous item \ts{11}{2}{11:08}\ts{11}{2}{11:29}\ts{11}{2}{11:52}.
  \item The Lie derivative acts on a function as the covariant derivative does: \textbf{true}, axiom 1, $\mathcal L_Xf=X(f)=\nabla_Xf$ \ts{11}{2}{12:14}.
\end{itemize}

\subsection*{Exercise 2: push-forward and pull-back formulae for practical use}

\begin{enumerate}
  \item \emph{The pull-back of a gradient is the gradient of the pull-back:} $\phi^*(\mathrm df)=\mathrm d(\phi^*f)$ for $f\in C^\infty(N)$. Act on an arbitrary $X\in T_pM$: $(\phi^*\mathrm df)(X)=\mathrm df(\phi_*X)=(\phi_*X)(f)=X(f\circ\phi)$, and the pull-back of a function is $\phi^*f:=f\circ\phi$ (the only thing one can get back on $M$ from $f$), so this is $X(\phi^*f)=\mathrm d(\phi^*f)(X)$; $X$ was arbitrary \ts{11}{2}{12:35}\ts{11}{2}{12:57}\ts{11}{2}{13:18}\ts{11}{2}{14:20}\ts{11}{2}{15:23}\ts{11}{2}{16:03}.
  \item \emph{Components of the push-forward:} $(\phi_*)^a{}_i=\mathrm dy^a\bigl(\phi_*\partial/\partial x^i|_p\bigr)=\frac{\partial}{\partial x^i}\big|_p(y^a\circ\phi)$, as derived in \cref{sec:c11-push} \ts{11}{2}{16:25}\ts{11}{2}{17:07}\ts{11}{2}{17:28}.
  \item \emph{Components of the pulled-back metric:} for $\phi\colon M\to N$ and $g$ on $N$ with components $g_{mn}$ in the chart $y$,
        \[
          (\phi^*g)_{ab}(p)=\frac{\partial(y^m\circ\phi)}{\partial x^a}\Big|_p\,\frac{\partial(y^n\circ\phi)}{\partial x^b}\Big|_p\,g_{mn}\bigl(\phi(p)\bigr).
        \]
        Proof: $(\phi^*g)_{ab}=g\bigl(\phi_*\partial_a,\phi_*\partial_b\bigr)$; expand both arguments in the chart basis $\partial/\partial y^m$ of $N$ at $\phi(p)$, the coefficients being the components of the push-forward (exercise 2.2); by bilinearity pull them out, leaving $g(\partial/\partial y^m,\partial/\partial y^n)=g_{mn}(\phi(p))$ \ts{11}{2}{17:48}\ts{11}{2}{18:31}\ts{11}{2}{19:14}\ts{11}{2}{20:05}\ts{11}{2}{21:30}\ts{11}{2}{22:15}\ts{11}{2}{24:00}\ts{11}{2}{24:44}.
\end{enumerate}

\subsection*{Exercise 3: the Lie derivative the pedestrian way}

\begin{enumerate}
  \item \emph{The embedding.} A smooth embedding $\eta\colon S^2\to\R^3$ is given in the charts $(\theta,\varphi)$ and the identity chart on $\R^3$, in terms of positive numbers $a,b,c$: presumably $\eta(\theta,\varphi)=(a\sin\theta\cos\varphi,\;b\sin\theta\sin\varphi,\;c\cos\theta)$ [reconstructed from the metric components found below; the sheet was not read]. What can one say about the shape of $\eta(S^2)$? \emph{Nothing}: $\R^3$ with the standard topology and atlas is just a smooth manifold; to talk about straight lines or shape one needs extra structure, a connection or a metric, and there is none \ts{11}{2}{25:05}\ts{11}{2}{25:29}\ts{11}{2}{26:12}\ts{11}{2}{26:32}\ts{11}{2}{26:53}.
  \item \emph{The ellipsoid metric.} Now equip $\R^3$ with the Euclidean metric $\delta_{mn}$ and compute the components of the induced metric $g^{\mathrm{ell}}=\eta^*\delta$ with the formula of exercise 2.3 (the tutor first writes the point in the wrong space, apologises and corrects: $\eta(\tilde p)=p$ with $\tilde p\in S^2$). The result is
        \begin{align*}
          g^{\mathrm{ell}}_{11}&=(a^2\cos^2\varphi+b^2\sin^2\varphi)\cos^2\theta+c^2\sin^2\theta,\\
          g^{\mathrm{ell}}_{12}=g^{\mathrm{ell}}_{21}&=(b^2-a^2)\sin\varphi\cos\varphi\sin\theta\cos\theta,\\
          g^{\mathrm{ell}}_{22}&=(a^2\sin^2\varphi+b^2\cos^2\varphi)\sin^2\theta,
        \end{align*}
        (the tutor leaves the middle component as $a^2\cos\varphi\cos\theta(-\sin\varphi\sin\theta)+b^2\sin\varphi\cos\theta\sin\theta\cos\varphi$, which is the same). The full calculation is in \cref{app:c11-ellipsoid}. For $a=b=c=1$ one gets the round metric of chapter~\ref{ch:c10}; the formula is long but easy to compute: from a metric on $\R^3$ and an embedding one pulls the metric back \ts{11}{2}{26:53}\ts{11}{2}{27:35}\ts{11}{2}{28:59}\ts{11}{2}{30:45}\ts{11}{2}{31:27}\ts{11}{2}{32:09}\ts{11}{2}{33:54}\ts{11}{2}{34:37}\ts{11}{2}{35:20}.
  \item \emph{The algebra.} Check that $X_1,X_2,X_3$ of \cref{ex:c11-so3} form a Lie subalgebra of $\Gamma(TS^2)$ and find the structure constants. Act on a test function $f$ so as not to forget terms. For $[X_1,X_3]f=X_1(X_3f)-X_3(X_1f)$: $X_1X_3f=-\sin\varphi\,\partial_\theta\partial_\varphi f-\cot\theta\cos\varphi\,\partial_\varphi^2f$; for the second term the derivative $\partial_\varphi$ acts on the \emph{coefficients} too (product rule), giving $-\cos\varphi\,\partial_\theta f-\sin\varphi\,\partial_\varphi\partial_\theta f+\cot\theta\sin\varphi\,\partial_\varphi f-\cot\theta\cos\varphi\,\partial_\varphi^2f$. In the difference the second-derivative terms cancel (``these two guys are the same with another sign''), leaving
        \[
          [X_1,X_3]f=\cos\varphi\,\partial_\theta f-\cot\theta\sin\varphi\,\partial_\varphi f=X_2f,
        \]
        so $[X_1,X_3]=1\cdot X_2$, structure constant $1$ (the tutor makes a few sign slips on the board and corrects them). The remaining pairs are a take-home exercise, ``and it works out very well'' \ts{11}{2}{35:43}\ts{11}{2}{36:04}\ts{11}{2}{36:46}\ts{11}{2}{37:28}\ts{11}{2}{38:52}\ts{11}{2}{40:20}\ts{11}{2}{41:46}\ts{11}{2}{42:31}\ts{11}{2}{43:15}.
  \item \emph{The integral curve of $X_3$ through $p=x^{-1}(\theta_0,\varphi_0)$.} It must satisfy $\gamma_p(0)=p$ and $v_{\gamma_p}=X_3$ along it. In the chart, $(\theta\circ\gamma)'\,\partial_\theta+(\varphi\circ\gamma)'\,\partial_\varphi=\partial_\varphi$, so $(\theta\circ\gamma)'=0$ and $(\varphi\circ\gamma)'=1$: integrating, $\theta(\gamma_p(\lambda))=\theta_0$ and $\varphi(\gamma_p(\lambda))=\lambda+\varphi_0$ (the constants fixed by the initial condition) \ts{11}{2}{43:35}\ts{11}{2}{44:16}\ts{11}{2}{45:00}\ts{11}{2}{46:06}\ts{11}{2}{46:49}\ts{11}{2}{47:34}\ts{11}{2}{48:19}.
  \item \emph{The flow and the Lie derivative.} The integral curves give the flow $h^{X_3}_\lambda$ of smooth maps $S^2\to S^2$; in the chart $\hat h_\lambda(\theta,\varphi)=(\theta,\varphi+\lambda)$: points on a parallel of latitude run around the sphere (the equator is one line, ``it is just cyclic''). Is $\mathcal L_{X_3}g^{\mathrm{ell}}=0$? Compute the pull-back of $g^{\mathrm{ell}}$ along the flow with exercise 2.3: since the derivative of $\theta$ with respect to $\theta$ is $1$ and the others are $0$ (the Jacobian is the identity, as only an offset in $\varphi$ appears), $(h_\lambda^*g^{\mathrm{ell}})_{ab}(\theta,\varphi)=g^{\mathrm{ell}}_{ab}(\theta,\varphi+\lambda)$. For the component $11$ this is $(a^2\cos^2(\varphi+\lambda)+b^2\sin^2(\varphi+\lambda))\cos^2\theta+c^2\sin^2\theta$, to be compared with $g^{\mathrm{ell}}_{11}$ at $\varphi$: they differ in the offset. \emph{If $a=b$} both equal $a^2\cos^2\theta+c^2\sin^2\theta$, and the same holds for all other components: then $h^{X_3*}_\lambda g^{\mathrm{ell}}=g^{\mathrm{ell}}$ for all $\lambda$, so $X_3$ is a symmetry and $\mathcal L_{X_3}g^{\mathrm{ell}}=0$: ``the ellipsoid deformed only in this direction'' is symmetric under rotations around the $c$-axis. This is the idea of the exercise: to see under which circumstances vector fields are symmetries \ts{11}{2}{49:07}\ts{11}{2}{49:28}\ts{11}{2}{50:50}\ts{11}{2}{51:33}\ts{11}{2}{52:16}\ts{11}{2}{54:24}\ts{11}{2}{55:47}\ts{11}{2}{56:29}\ts{11}{2}{57:13}\ts{11}{2}{58:16}\ts{11}{2}{59:20}.
\end{enumerate}
\begin{coursenote}[Added by the notes: the converse and the shortcut]
(Not said in class.) If $a\neq b$ the $11$ component depends on $\varphi$, so $X_3$ is \emph{not} a symmetry: ``if and only if $a=b$''. The same follows without any flow from the component formula of the Lie derivative, since $\partial_iX_3^s=0$ in this chart: $(\mathcal L_{X_3}g)_{ij}=\partial_\varphi g_{ij}$, which vanishes for $g^{\mathrm{ell}}$ exactly if its components do not depend on $\varphi$, i.e.\ if $a=b$.
\end{coursenote}

\section*{Key concepts}
\begin{keyconcepts}
\textbf{Retrieval questions.}
\begin{enumerate}
  \item Define the push-forward $\phi_*$ and the pull-back $\phi^*$ of a smooth map: what do they act on, and what are their components? Why is a vector field not generally pushed forward to a vector field?
  \item Prove that $\phi_*$ sends the tangent vector of a curve to the tangent vector of the pushed-forward curve.
  \item How does one define the induced metric on an embedded manifold? What is the induced metric of $S^2\hookrightarrow\R^3$ for the embedding with axes $a,b,c$?
  \item What is an integral curve? When is a vector field complete? Give an example of an incomplete field and say why this matters for singularities. What is the flow $h^X_\lambda$?
  \item What is a Lie algebra, a Lie subalgebra and its structure constants? Give the rotation fields on $S^2$.
  \item Define: a Lie subalgebra is a symmetry of a metric; an isometry. Why is the $a=b\neq c$ ellipsoid not $SO(3)$-symmetric?
  \item State the axioms of the Lie derivative and why it needs no extra structure. Write $(\mathcal L_Xg)_{ij}$ in a chart and use it to test $X_3=\partial_\varphi$ on the ellipsoid.
\end{enumerate}
\textbf{Concepts.}
\begin{itemize}
  \item Push-forward and pull-back, components $\partial\hat\phi^a/\partial x^i$, curves and level lines (\cref{def:c11-push}, \cref{thm:c11-curve}, \cref{def:c11-pull}).
  \item Induced metric $g_M=\phi^*g_N$ (\cref{def:c11-induced}).
  \item Integral curve, complete field, flow (\cref{def:c11-integral}, \cref{def:c11-complete}, \cref{def:c11-flow}).
  \item Lie algebra of vector fields, Lie subalgebra, structure constants, $\mathfrak{so}(3)$ on $S^2$ (\cref{def:c11-liealg}, \cref{def:c11-subalg}, \cref{ex:c11-so3}).
  \item Symmetry of a metric, isometry (\cref{def:c11-symmetry}); Lie derivative $\mathcal L_Xg=0$ (\cref{def:c11-liederivative}).
\end{itemize}
\end{keyconcepts}

\flushside
\end{paracol}
